\documentclass[11pt]{article}
\usepackage[OT1]{fontenc}
\renewcommand{\rmdefault}{cmr}
\renewcommand{\sfdefault}{cmss}
\renewcommand{\ttdefault}{cmtt}
\usepackage{amsmath,amssymb,amsthm,mathtools}
\usepackage[margin=27mm]{geometry}
\usepackage{microtype,booktabs,longtable,array}
\usepackage[hidelinks]{hyperref}
\setlength{\emergencystretch}{2em}
\setlength{\parskip}{3pt plus 1pt}
\setcounter{tocdepth}{1}
\newtheorem{theorem}{Theorem}[section]
\newtheorem{lemma}[theorem]{Lemma}
\newtheorem{proposition}[theorem]{Proposition}
\newtheorem{corollary}[theorem]{Corollary}
\theoremstyle{definition}
\newtheorem{definition}[theorem]{Definition}
\newtheorem{example}[theorem]{Example}
\theoremstyle{remark}
\newtheorem{remark}[theorem]{Remark}
\newcommand{\N}{\mathbb N}
\newcommand{\Q}{\mathbb Q}
\newcommand{\R}{\mathbb R}
\newcommand{\Z}{\mathbb Z}
\newcommand{\supp}{\operatorname{supp}}
\newcommand{\Con}{\operatorname{Con}}
\newcommand{\ord}{\operatorname{ord}}
\newcommand{\dom}{\operatorname{dom}}
\newcommand{\sgn}{\operatorname{sgn}}
\newcommand{\res}{\operatorname{Res}}
\newcommand{\rad}{\operatorname{rad}}
\newcommand{\CK}{\mathsf{CK}}
\newcommand{\KT}{\mathsf{KT}}
\newcommand{\DCB}{\mathsf{DCB}}
\newcommand{\FS}{\mathsf{FS}}
\newcommand{\Ag}{\mathsf{Ag}}
\newcommand{\Know}{\mathsf K}
\newcommand{\Truth}{\mathsf T}
\newcommand{\Prv}{\mathsf{Pr}}
\newcommand{\code}[1]{\ulcorner #1\urcorner}
\newcommand{\Tot}{\operatorname{Tot}}
\numberwithin{equation}{section}

\title{Explicit Constraint Decisions for Differential Polynomials\\over the Rational Constants}
\author{}
\date{28 October 2026\\Working manuscript}
\begin{document}
\maketitle
\begin{abstract}
The constrained-pair predicate asks whether every two solutions of a
differential equation, after excluding a lower-order zero set, generate
isomorphic differential fields with the designated solutions identified.
We give exact decisions for two classes over the rational field with
zero derivation. For an irreducible first-order polynomial $p(Y,DY)$,
put $r(Y)=p(Y,0)$. If $r\ne0$, a nonzero polynomial $q(Y)$ is a
constraint exactly when the squarefree part of $r$ divides $q$; if
$r=0$, none is. For the linear equation $f(D)Y=0$, where the roots of
the monic rational polynomial $f$ are linearly independent over $\Q$,
the constraint criterion is divisibility by an explicit rational
product of spectral linear forms. A resultant computes this product
without manipulating the roots. Cyclotomic examples give an elementary
recognizable family of unbounded differential orders. We prove both
directions of each criterion and construct the needed solutions
algebraically. These are restricted decisions for the actual
constrained-pair predicate, not a decision procedure for all
differential polynomials over $\Q$.
\end{abstract}
\tableofcontents

\section{A differential analogue of irreducibility}

For an ordinary irreducible polynomial over $\Q$, any two roots
generate fields isomorphic over $\Q$ by an isomorphism sending one
root to the other. A differential equation is more complicated.
Two solutions of the same equation may satisfy different additional
differential equations and need not generate the same kind of
differential field. A constraint excludes lower-order exceptional
solutions until a single designated-root isomorphism type remains.

Miller, Ovchinnikov and Trushin \cite{MOT} study the computational
role of these pairs as an analogue of irreducible polynomials.
Their broad question includes the computability of constraint
sets. Even with base $\Q$ and zero derivation, a general procedure
is not supplied by their results. Here we treat two concrete input
classes. In each case the problem reduces to ordinary rational
polynomial divisibility.

The two mechanisms are different. In order one, all nonconstant
solutions of an irreducible equation have the same differential
function field; one must exclude its constant solutions. In higher
order, spectral independence forces all nonzero spectral coordinates
to be algebraically independent; one must exclude solutions missing
a coordinate. This distinction also explains why the first-order
argument cannot simply be iterated.

\section{Definitions and algebraic preliminaries}

A differential field is a field $F$ with a derivation $D$, satisfying
$D(a+b)=Da+Db$ and $D(ab)=aDb+bDa$. All fields here have characteristic
zero, there is one derivation, and $\Q$ has its zero derivation.
Write
\[
 \Q\{Y\}=\Q[Y_0,Y_1,Y_2,\ldots],\qquad DY_i=Y_{i+1}.
\]
Evaluation at $a$ sends $Y_i$ to $D^ia$. The order of a nonconstant
differential polynomial is the largest derivative index it uses;
nonzero constants are treated as having order below zero. Write
$\Q\langle a\rangle$ for the differential field generated by $a$.

\begin{definition}\label{def:constraint}
Let $p$ be a source-normalized, algebraically irreducible
differential polynomial of order $n$, and let $q\ne0$ have order
less than $n$. The pair $(p,q)$ is \emph{constrained over $\Q$} if
there is a solution of $p(a)=0$, $q(a)\ne0$ in a differential
extension and any two such solutions $a,b$ in arbitrary differential
extensions admit a differential-field isomorphism
\[
 \Q\langle a\rangle\longrightarrow\Q\langle b\rangle,
 \qquad a\longmapsto b,
\]
fixing $\Q$. Denote the positive predicate by $\Con_\Q(p,q)$.
\end{definition}

This is the designated-root condition of \cite[Definition~4.3]{MOT}.
Their notation $T_\Q$ is the complementary predicate, so signs must
be kept distinct when translating statements. Scalar normalization
of $p$ does not change its solutions. In the first-order theorem
we assume that source normalization, not the stronger condition that
the coefficient of the highest derivative be 1.

\begin{lemma}\label{lem:constants}
An element algebraic over $\Q$ in a differential extension is a
constant. A nonconstant element is therefore transcendental over $\Q$.
\end{lemma}
\begin{proof}
If $m(a)=0$ is its irreducible rational polynomial, differentiation
gives $m'(a)Da=0$. Separability in characteristic zero gives
$m'(a)\ne0$, hence $Da=0$.
\end{proof}

We also use the elementary fact that a derivation on a
characteristic-zero field extends uniquely to an algebraic extension.
For an algebraic element $a$ with minimal polynomial
$\sum_i c_iT^i$, its required derivative is
\[
 Da=-\frac{\sum_i(Dc_i)a^i}{\sum_i i c_i a^{i-1}}.
\]
The denominator is nonzero by separability. This formula extends the
derivation successively to finite algebraic extensions and then to
an algebraic closure; uniqueness makes the extensions compatible.
In particular every algebraic number remains constant.

\section{Every first-order input in the stated slice}

Let $p(Y,Z)\in\Q[Y,Z]$ be algebraically irreducible and have positive
degree in $Z$, representing $Z=DY$. Let $0\ne q(Y)\in\Q[Y]$, and
put $r(Y)=p(Y,0)$. For $r\ne0$, let $\rad(r)$ be its monic squarefree
part, with $\rad(r)=1$ for nonzero constant $r$.

\begin{theorem}\label{thm:orderone}
On these source-valid inputs:
\begin{enumerate}
\item if $r\ne0$, then $\Con_\Q(p,q)$ holds exactly when
      $\rad(r)\mid q$ in $\Q[Y]$;
\item if $r=0$, then $\Con_\Q(p,q)$ fails for every nonzero $q$.
\end{enumerate}
\end{theorem}

\subsection{The common field of all nonconstant solutions}

\begin{lemma}\label{lem:curve}
Any two nonconstant solutions of $p(a,Da)=0$ generate differential
fields isomorphic over $\Q$ by a map sending one solution to the other.
\end{lemma}
\begin{proof}
Such an $a$ is transcendental by Lemma~\ref{lem:constants}.
The polynomial $p$ is primitive as a polynomial in $Z$ over $\Q[Y]$:
otherwise its nonconstant content would factor it. Gauss's lemma
makes it irreducible over $\Q(Y)$. Substitution $Y\mapsto a$
therefore makes $p(a,Z)$ the minimal polynomial of $Da$, up to a
nonzero scalar over $\Q(a)$. It follows that the evaluation kernel
of $\Q[Y,Z]\to\Q[a,Da]$ is exactly $(p)$. One can see the contraction
directly: an element of $\Q[Y,Z]$ divisible by primitive $p$ in
$\Q(Y)[Z]$ is divisible in $\Q[Y,Z]$.

The partial derivative $p_Z$ is nonzero and has smaller $Z$-degree
than $p$. It is not in $(p)$, so $p_Z(a,Da)\ne0$. Differentiating
the equation gives
\begin{equation}\label{eq:second}
 D^2a=-\frac{p_Y(a,Da)Da}{p_Z(a,Da)}.
\end{equation}
Thus $F_a=\Q(a,Da)$ is closed under the derivation, and equals
$\Q\langle a\rangle$. For another nonconstant solution $b$,
the two evaluation maps identify both fields with
$\operatorname{Frac}(\Q[Y,Z]/(p))$. The resulting isomorphism sends
$a,Da$ to $b,Db$. It commutes with $D$ on the first generator by
construction, on the second by \eqref{eq:second}, and hence on the
whole field by the derivation rules.
\end{proof}

The separant $p_Z(a,Da)$ is therefore nonzero at every nonconstant
solution in this setting. It is a conclusion of irreducibility and
transcendence, not a hidden genericity assumption.

\begin{lemma}\label{lem:genericcurve}
If $r\ne0$, there exists a nonconstant solution avoiding every
specified nonzero $q(Y)$.
\end{lemma}
\begin{proof}
In the domain $\Q[Y,Z]/(p)$ let $y,z$ be the residue classes, and
let $F$ be its fraction field. The class $y$ is transcendental over
$\Q$, and $p_Z(y,z)\ne0$. Set
\[
 Dy=z,\qquad Dz=-p_Y(y,z)z/p_Z(y,z).
\]
These prescriptions respect the defining relation because
$D(p(y,z))=0$, so they define a derivation on the fraction field.
Equivalently extend the usual formal derivation on the localized
polynomial ring and pass through the relation. Since $r\ne0$,
irreducible $p$ is not a scalar multiple of $Z$, so $z\ne0$.
Thus $y$ is nonconstant, solves $p(y,Dy)=0$, and avoids $q$ by
transcendence. This is an explicit algebraic differential extension
of $\Q$; no analytic solution or differentially closed field is needed.
\end{proof}

\subsection{Removing precisely the constant exceptions}

\begin{proof}[Proof of Theorem~\ref{thm:orderone}]
Assume $r\ne0$. Constant solutions are exactly the algebraic constants
with $r(a)=0$. Indeed a constant solution satisfies that equation;
conversely every root of $r$ in an algebraic closure of $\Q$, equipped
with the zero derivation, is a solution. A nonzero rational polynomial
vanishes at all these roots precisely when it is divisible by
$\rad(r)$.

If $\rad(r)\mid q$, any solution avoiding $q$ is therefore
nonconstant. Lemma~\ref{lem:curve} gives the required isomorphism
between any two, and Lemma~\ref{lem:genericcurve} supplies existence.

If $\rad(r)\nmid q$, take an irreducible factor of $r$ not dividing
$q$ and one of its algebraic roots $c$. Then $q(c)\ne0$ and $c$
is a constant solution. The nonconstant solution from
Lemma~\ref{lem:genericcurve} also avoids $q$. A differential isomorphism
sending one to the other would preserve vanishing of $DY$, which is
impossible. Hence the pair is not constrained.

If $r=0$, then $Z$ divides $p$. Irreducibility forces $p$ to be a
nonzero rational multiple of $Z$. Infinitely many rational constants
avoid $q$. Two distinct such constants are solutions but cannot be
mapped to each other by an isomorphism fixing $\Q$. This proves
the second case.
\end{proof}

\begin{corollary}\label{cor:algorithmone}
The constrained-pair predicate is decidable on this first-order slice.
\end{corollary}
\begin{proof}
Compute $r$ by substitution. If zero, answer no. Otherwise compute
$\gcd(r,r')$, form the monic associate of $r/\gcd(r,r')$, and test
its divisibility into $q$. These are finite rational polynomial
operations and implement the theorem.
\end{proof}

\begin{example}
For $p=Z-Y$, $r=-Y$, so exactly the nonzero $q$ divisible by $Y$
are constraints. The equation has the exceptional solution zero;
every nonzero solution has the same designated-root differential field.
For $p=YZ-1$, $r=-1$, so every nonzero $q(Y)$ works. For $p=Z$,
none works, since lower-order exclusion by one nonzero polynomial
cannot reduce all rational constants to a single designated value.
\end{example}

At order zero, source-valid irreducible $p(Y)$ and nonzero constant
$q$ are constrained for the usual algebraic reason: all roots are
algebraic constants and their ordinary $\Q$-isomorphisms are
differential isomorphisms. This small case should not be confused
with the first-order equation $DY=0$.

\section{Higher-order equations with independent spectral values}

Let
\[
 f(T)=T^n+a_{n-1}T^{n-1}+\cdots+a_0\in\Q[T],\qquad n\ge1,
\]
and suppose its roots $\lambda_1,\ldots,\lambda_n$ in a splitting
field $K$ are linearly independent over $\Q$. In particular they
are nonzero and pairwise distinct. Put
\[
 p(Y)=f(D)Y=Y_n+a_{n-1}Y_{n-1}+\cdots+a_0Y_0.
\]
This differential polynomial is algebraically irreducible because
it is monic linear in $Y_n$, whether or not $f$ is irreducible as
an ordinary univariate polynomial. Let
$0\ne q\in\Q[Y_0,\ldots,Y_{n-1}]$.

For each $i$, define the linear form
\[
 L_i(Y)=\left[\prod_{j\ne i}(D-\lambda_j)\right]Y,
 \qquad R_f(Y)=\prod_{i=1}^n L_i(Y).
\]
The Galois group permutes these forms, so their product has rational
coefficients. It uses only the first $n$ jets and is nonzero.

\begin{theorem}\label{thm:spectral}
Under this spectral independence hypothesis,
\[
 \Con_\Q(p,q)\quad\Longleftrightarrow\quad
 R_f\mid q\quad\text{in }\Q[Y_0,\ldots,Y_{n-1}].
\]
In particular $(p,R_f)$ is a positive constrained pair.
\end{theorem}

\subsection{Spectral coordinates are valid for every solution}

Let $a$ be any solution in any differential extension. Extend its
field algebraically so that it contains a copy of $K$. Its elements
are constants by Lemma~\ref{lem:constants}. Define
\[
 x_i=L_i(a)/f'(\lambda_i).
\]
Lagrange interpolation gives, for $0\le k<n$,
\begin{equation}\label{eq:vandermonde}
 D^ka=\sum_i\lambda_i^k x_i,
 \qquad Dx_i=\lambda_i x_i.
\end{equation}
For the first identity, interpolate $T^k$ at the distinct roots
using the polynomials
$\prod_{j\ne i}(T-\lambda_j)/f'(\lambda_i)$, then apply the resulting
operator identity to $a$. For the second use
$(D-\lambda_i)L_i(a)=p(a)=0$.
The Vandermonde matrix in \eqref{eq:vandermonde} is invertible, so
these are invertible linear coordinates on the jets, and
\begin{equation}\label{eq:product}
 R_f(a)=\left(\prod_i f'(\lambda_i)\right)\prod_i x_i.
\end{equation}

\begin{lemma}\label{lem:weights}
If all $x_i$ are nonzero and $Dx_i=\lambda_ix_i$, they are
algebraically independent over $K$.
\end{lemma}
\begin{proof}
For a multi-index $\alpha\in\N^n$, the monomial $x^\alpha$ has
derivative $w_\alpha x^\alpha$, with
$w_\alpha=\sum_i\alpha_i\lambda_i$. Rational linear independence
of the $\lambda_i$ makes these weights distinct for distinct
multi-indices.

Suppose a nonzero polynomial relation
$\sum_{\alpha\in S}c_\alpha x^\alpha=0$ holds, after combining
equal monomials and discarding zero coefficients. Fix $\alpha_0\in S$
and apply the differential operator
$\prod_{\beta\in S\setminus\{\alpha_0\}}(D-w_\beta)$.
The coefficients and weights are constants. Every term except
$\alpha_0$ vanishes, leaving
\[
 c_{\alpha_0}\prod_{\beta\ne\alpha_0}
 (w_{\alpha_0}-w_\beta)x^{\alpha_0}=0.
\]
All factors are nonzero in a field, a contradiction. This also
handles a constant monomial if the relation has one.
\end{proof}

\begin{lemma}\label{lem:genericjets}
If $p(a)=0$ and $R_f(a)\ne0$, its first $n$ jets are algebraically
independent over $\Q$. Any two such solutions have the required
designated-root differential isomorphism over $\Q$.
\end{lemma}
\begin{proof}
By \eqref{eq:product}, all spectral coordinates are nonzero.
Lemma~\ref{lem:weights} and the invertible change of coordinates
give jet independence even over $K$. The recurrence $p(a)=0$
expresses $D^na$ as the fixed rational linear combination of lower
jets. Thus
$F_a=\Q(a,Da,\ldots,D^{n-1}a)$ is closed under $D$ and equals
$\Q\langle a\rangle$. For a second solution $b$, algebraic
independence gives an isomorphism sending each $D^ka$ to $D^kb$.
It commutes with the derivation on the first $n-1$ generators by
construction and on the last by the common recurrence. Hence it
is a differential isomorphism on the whole field, fixing $\Q$
and sending $a$ to $b$.
\end{proof}

\subsection{Generic and missing-coordinate solutions}

Use algebraically independent variables $X_1,\ldots,X_n$ over $K$
and give the rational function field $K(X_1,\ldots,X_n)$ the
derivation $DX_i=\lambda_iX_i$, zero on $K$. This is a valid
differential field by the polynomial and quotient rules. The element
$b=\sum_iX_i$ solves $p(b)=0$. Its jets are the invertible Vandermonde
transform of the $X_i$, so every nonzero rational polynomial $q$
in those jets is nonzero at $b$. Also $R_f(b)\ne0$. This supplies
the existence needed in Definition~\ref{def:constraint}.

If one sets $X_i=0$ and keeps the other variables independent, the
smaller rational function field $K(X_j:j\ne i)$, with the same
derivation rule, gives a solution $c=\sum_{j\ne i}X_j$ with a
missing coordinate. It satisfies $R_f(c)=0$.

\begin{proof}[Proof of Theorem~\ref{thm:spectral}]
If $R_f\mid q$, a solution avoiding $q$ has $R_f\ne0$.
Lemma~\ref{lem:genericjets} gives the desired isomorphism between
any two, and the all-coordinate construction supplies existence.

Suppose $R_f\nmid q$. Divisibility does not change on extending
coefficients to $K$: if $q=R_fh$ over $K$, Galois invariance and
uniqueness of the quotient imply $h\in\Q[Y_0,\ldots,Y_{n-1}]$.
The $L_i$ are pairwise nonassociate linear coordinate forms over $K$.
Since their product does not divide $q$, some $L_i$ does not divide
$q$. In spectral coordinates, the restriction of $q$ to $X_i=0$
is therefore nonzero.

Choose the missing-coordinate solution $c$ just constructed. Its
remaining coordinates are independent, so $q(c)\ne0$, while
$R_f(c)=0$. Compare it with the all-coordinate solution $b$, for
which both $q(b)$ and $R_f(b)$ are nonzero. The differential
polynomial $R_f$ over $\Q$ distinguishes them, so no
designated-root differential isomorphism exists. For $n=1$ the
missing-coordinate field is simply $K$ and $c=0$; the same
argument applies. This proves necessity.
\end{proof}

\section{A rational algorithm and explicit unbounded examples}

The proof used a splitting field, but the decision polynomial can
be computed over $\Q$ alone. Write $a_n=1$ and set
\[
 G_f(T)=\sum_{k=1}^n a_k\sum_{j=0}^{k-1}T^jY_{k-1-j}.
\]
It is obtained from $(f(S)-f(T))/(S-T)$ by replacing each $S^m$
with $Y_m$. Substituting $T=\lambda_i$ gives $L_i$, so
\begin{equation}\label{eq:resultant}
 R_f=\res_T(f(T),G_f(T)).
\end{equation}
Here the resultant identity uses that $f$ is monic. A Sylvester
determinant computes this polynomial in rational arithmetic.
Multivariate polynomial division by the single polynomial $R_f$
then decides whether $R_f\mid q$.

This is a terminating decision algorithm on the spectral promise.
It does not need the roots as additional input. No implementation
of a general promise recognizer is asserted. The following
subclass has an elementary recognizable promise.

\begin{example}[A quadratic equation]
Take $f(T)=T^2-T-1$. Its roots
$(1+\sqrt5)/2,(1-\sqrt5)/2$ are rationally independent: the rational
and $\sqrt5$ parts of a rational relation force both coefficients
to vanish. Then
\[
 p=Y_2-Y_1-Y_0,\qquad
 R_f=Y_1^2-Y_0Y_1-Y_0^2.
\]
Every nonzero $q(Y_0,Y_1)$ is decided: the pair is constrained
exactly when this displayed quadratic divides $q$. In particular
$q=R_f$ gives a positive pair and $q=1$ a negative one.
\end{example}

\begin{proposition}\label{prop:cyclotomic}
For every prime $\ell$, the polynomial
$f(T)=1+T+\cdots+T^{\ell-1}$ satisfies the spectral promise.
Consequently the theorem supplies positive pairs and full decisions
for all nonzero $q$ at an unbounded set of differential orders.
\end{proposition}
\begin{proof}
The roots are the nontrivial $\ell$th roots
$\zeta,\ldots,\zeta^{\ell-1}$. The polynomial $f$ is irreducible:
substitute $T=U+1$ and apply Eisenstein at $\ell$ to
$((U+1)^\ell-1)/U$. If a rational relation among the displayed
roots existed, it would be a polynomial of degree at most
$\ell-1$ with zero constant term vanishing at $\zeta$.
Its minimal polynomial is $f$, so this relation polynomial would
be a scalar multiple of $f$. The zero constant term forces that
scalar to be zero. The roots are therefore rationally independent.
The orders are $\ell-1$, unbounded as primes grow. Trial division
and checking the coefficients recognize this family.
\end{proof}

\section{Why dependent spectra require a different argument}

Spectral independence is not a disposable sufficient condition.
If the weights are dependent, nonzero coordinates can carry
additional rational invariants, distinguishing solutions which
all avoid the proposed product.

For $f(T)=T^2-1$, the roots $1,-1$ are dependent, and the analogous
product is $R=Y_1^2-Y_0^2$. In $\Q(x)$ with $Dx=x$, put
\[
 a_c=(x+c/x)/2\qquad(c\in\Q\setminus\{0\}).
\]
Then $D^2a_c=a_c$ and $R(a_c)=-c$. The choices $c=1,2$ both
avoid $R$ but are distinguished by $R+1$, which vanishes at only
the first. Thus even $(Y_2-Y_0,R)$ is not constrained. The
independent-spectrum theorem would give a false answer here
if its hypothesis were omitted.

At order one the exceptional solutions were precisely the
constants. At higher order a solution may have several nonzero
derivatives and still satisfy additional algebraic or differential
relations among its jets. The spectral proof succeeds by proving
jet independence, not merely by observing nonconstancy.

\section{Relation to the general constraint-set problem}

The first theorem decides every supplied nonzero lower-order $q$
for the specified irreducible order-one inputs over the pure rational
constant field. The second decides every supplied nonzero $q$ for
an unbounded-order family. These are decisions for the designated-root
predicate itself, including both existence and all-extension
isomorphism requirements. They do not replace $\Q$ by a field
with a named nonconstant parameter or ask only whether some
constraint exists.

The proofs use standard differential-field and polynomial methods;
the local contribution is the explicit criterion and its effective
form for these inputs. The general predicate for repeated roots,
dependent spectra and nonlinear higher-order equations is not
decided here. Neither theorem establishes the broader differential
Kronecker or extension-property conclusions discussed in \cite{MOT}.

A bounded source and literature check did not identify an equivalent
published formulation of both criteria. This is not a determination
of historical novelty; the first-order function-field observation
and the spectral independence argument in particular belong to
classical differential algebra. Further attribution work is
appropriate before claiming these as new literature results.

\paragraph{Preparation.}
This AI-assisted hand-proof manuscript includes all local algebraic
constructions and both directions of the decision criteria. The
algorithms are proved terminating on their stated input promises;
no empirical verification or formal certification is claimed.

\begin{thebibliography}{9}
\bibitem{MOT} R.~Miller, A.~Ovchinnikov and D.~Trushin,
\emph{Computing Constraint Sets for Differential Fields},
Journal of Algebra 407 (2014), 316--357.
\href{https://doi.org/10.1016/j.jalgebra.2014.02.032}{doi:10.1016/j.jalgebra.2014.02.032}.
Source version \href{https://arxiv.org/abs/1208.1152v2}{arXiv:1208.1152v2},
12 April 2014; especially Definition~4.3 and the constraint-set
decidability discussion.
\end{thebibliography}
\end{document}
